\chapter{Lezione 1}

\section{Introduzione al Corso e alla Formazione delle Immagini}
Il corso di Computer Vision inizia affrontando il problema classico dell'acquisizione e della formazione delle immagini.Trattandosi di una disciplina incentrata sull'elaborazione visiva tramite elaboratore, è indispensabile comprendere prima di tutto i meccanismi fisici e geometrici attraverso cui si generano le immagini.

La formazione delle immagini è un processo fortemente multidisciplinare che coinvolge:
\begin{itemize}
    \item \textbf{Matematica}: in particolare la geometria dello spazio e delle proiezioni.
    \item \textbf{Fisica e Ottica}: studio della propagazione della luce e del comportamento delle lenti.
    \item \textbf{Ingegneria Elettronica}: funzionamento dei sensori e conversione del segnale luminoso in potenziale elettrico.
\end{itemize}

Nell'ambito della Computer Vision non è necessario approfondire eccessivamente la fisica quantistica o l'ingegneria elettronica dei sensori, ma è essenziale possedere un modello geometrico e concettuale chiaro di come il mondo reale tridimensionale si trasformi in un'immagine bidimensionale. L'obiettivo primario del corso è infatti imparare ad rielaborare ed estrarre informazioni geometriche e strutturali dalla scena a partire dalle immagini acquisite.

Lo strumento di acquisizione universale è la macchina fotografica (o fotocamera digitale, come quelle integrate nei moderni smartphone).

\section{Componenti Principali di un Sistema di Acquisizione Ottico}
Un generico sistema di acquisizione ottica si compone dei seguenti elementi geometrici e fisici fondamentali:
\begin{itemize}
    \item \textbf{Piano Immagine / Schermo}: È l'elemento reale e fisico su cui si forma l'immagine. Può consistere in una pellicola chimica (fotografia analogica) o in un sensore elettronico CCD/CMOS (fotografia digitale).
    \item \textbf{Asse Ottico}: È un'entità teorica e geometrica ortogonale al piano dello schermo. Definisce l'orientamento della telecamera rispetto alla scena tridimensionale.
    \item \textbf{Lente / Sistema Ottico}: Posizionata tra la scena e il piano immagine, è generalmente costituita da un insieme complesso di lenti in vetro progettate per focalizzare la luce in diverse condizioni operative.
    \item \textbf{Apertura / Diaframma}: Un meccanismo di regolazione del passaggio della luce. Senza un vincolo sull'apertura, l'eccesso di luce sovraesporrebbe l'immagine (rendendola "completamente bianca"). L'ampiezza dell'apertura è strettamente bilanciata con il tempo di scatto o tempo di esposizione.
\end{itemize}

\section{Geometria della Proiezione: Riduzione di Dimensione $3\text{D} \to 2\text{D}$}
Il processo di formazione dell'immagine si basa su una proiezione geometrica: la luce colpisce gli oggetti della scena $3\text{D}$, viene riflessa e proiettata sullo schermo $2\text{D}$.

\begin{equation}
    \text{Scena } 3\text{D} \; (X, Y, Z) \longrightarrow \text{Immagine } 2\text{D} \; (u, v)
\end{equation}

In questa trasformazione si verifica una \textbf{perdita irreversibile di informazione}, ovvero la coordinata di profondità $Z$. La ricostruzione della profondità da un'unica immagine è matematicamente mal posta; la percezione della terza dimensione nell'essere umano avviene grazie alla conoscenza a priori della struttura del mondo reale e non direttamente dalla singola proiezione.

I quattro aspetti principali coinvolti nella formazione delle immagini sono:
\begin{enumerate}
    \item \textbf{Geometria}: Modellazione geometrica dei raggi e dei piani di proiezione.
    \item \textbf{Ottica}: Comportamento e deviazione della luce attraverso i vetri della lente.
    \item \textbf{Radiometria}: Misurazione della quantità di energia luminosa irradiata e riflessa.
    \item \textbf{Sensori}: Conversione dell'energia luminosa in carica o differenza di potenziale elettrico.
\end{enumerate}
Nel corso ci si focalizzerà quasi esclusivamente sulla \textbf{Geometria}, costruendo modelli matematici e sintetici di telecamera.

\section{Tipologie di Proiezione: Ortografica vs. Prospettica}
Confrontando due immagini reali, ad esempio la foto del Pentagono dall'alto e la foto di un ponte visto dal basso, si notano differenze geometriche sostanziali:

\subsection{Proiezione Ortografica}
Nella proiezione ortografica, i raggi di proiezione sono tutti paralleli tra loro e ortogonali al piano immagine. 
\begin{itemize}
    \item Conserva le dimensioni e le proporzioni reali degli oggetti senza distorsione prospettica.
    \item È un'approssimazione valida quando la distanza dell'osservatore dalla scena è molto grande rispetto all'estensione/profondità dell'oggetto stesso (es. vista aerea o satellitare di un edificio di $10\text{ m}$ acquisita da $2\text{--}3\text{ km}$ di altezza).
    \item Le equazioni di proiezione semplificate sono:
    \begin{equation}
        u = X, \quad v = Y
    \end{equation}
    poiché le variazioni di profondità $\Delta Z$ sono considerate trascurabili.
\end{itemize}

\subsection{Proiezione Prospettica}
È il modello più naturale e realistico, corrispondente alla visione umana e alla resa pittorica (sviluppata dal Rinascimento e favorita dall'uso della camera oscura/stenopeica).
\begin{itemize}
    \item I raggi di proiezione non sono paralleli, ma convergono tutti in un singolo punto dello spazio detto \textbf{Centro Ottico} o \textbf{Centro di Proiezione} ($C$).
    \item Gli oggetti paralleli e le linee parallele nello spazio convergono visivamente verso un punto dell'immagine.
    \item Introduce una proporzionalità inversa rispetto alla profondità $Z$: man mano che un oggetto si allontana, la sua proiezione sull'immagine diminuisce di dimensione.
\end{itemize}

\section{Il Modello Geometrico Pinhole (Punto di Spillo)}
Il modello \textit{Pinhole} (o foro stenopeico / punta di spillo) rappresenta il modello matematico ideale di telecamera.

\subsection{Assunzioni del Modello}
Si assume che l'apertura della telecamera sia ridotta a un singolo **punto geometrico** $C$ (Centro di Proiezione).
\begin{itemize}
    \item Per ogni punto della scena $3\text{D}$ passa **un unico raggio ottico** che attraversa il centro di proiezione $C$ e colpisce il piano immagine.
    \item Questo modello astratto elimina ogni problema di aberrazione o sfocatura: **tutti i punti della scena sono sempre perfettamente a fuoco**, indipendentemente dalla loro distanza $Z$.
\end{itemize}

\subsection{Geometria e Triangoli Simili}
Fissiamo un sistema di riferimento tridimensionale cartesiano centrato in $C$:
\begin{itemize}
    \item Origine $(0,0,0)$ coincide con il Centro di Proiezione $C$.
    \item L'asse $Z$ coincide con l'Asse Ottico ed è ortogonale al piano immagine.
    \item Gli assi $X$ e $Y$ sono paralleli al piano immagine.
    \item Il piano immagine si trova a una distanza $f$ (distanza focale) dal punto $C$ lungo l'asse $Z$.
\end{itemize}

Considerando un punto nello spazio $P = (X, Y, Z)^T$ e la sua proiezione $p = (u, v)^T$ sul piano immagine, si formano due triangoli rettangoli simili opposti al vertice $C$. Sfruttando la proporzionalità tra i lati dei triangoli simili si ricava:

\begin{equation}
    \frac{u}{X} = -\frac{f}{Z} \implies u = -f \frac{X}{Z}
\end{equation}
\begin{equation}
    \frac{v}{Y} = -\frac{f}{Z} \implies v = -f \frac{Y}{Z}
\end{equation}

Il segno meno ($-$) indica matematicamente che l'immagine proiettata sul piano reale dietro al centro di proiezione è **capovolta** sia in senso verticale che orizzontale.

\subsection{Modello Virtuale Equivalentemente Non Invertito}
Per comodità algebrica ed evitare di gestire il segno meno, si definisce spesso un piano immagine virtuale posto \textbf{davanti} al centro di proiezione ad una distanza $+f$. Geometria e regole rimangono identiche, ma le equazioni di proiezione diventano:
\begin{equation}
    u = f \frac{X}{Z}, \quad v = f \frac{Y}{Z}
\end{equation}

\subsection{Modello Debolmente Prospettico (Weak Perspective)}
Quando la variazione di profondità dei punti della scena $\Delta Z$ è molto piccola rispetto alla distanza media $Z_0$ dall'osservatore, si può approssimare $Z \approx Z_0 = \text{costante}$. La proiezione equivale a una proiezione ortografica seguita da un ridimensionamento globale (scaling) in funzione della profondità media $Z_0$.

\section{Sistemi di Riferimento: Mondo vs. Telecamera e Parametri}
In uno scenario reale, non è possibile misurare la posizione dei punti $3\text{D}$ direttamente rispetto al sistema di riferimento interno della telecamera (poiché l'origine $C$ si trova all'interno dell'ottica). 

Occorre quindi definire un **Sistema di Riferimento Mondo** esterno arbitrario. La trasformazione dal Sistema Mondo al Sistema Telecamera è una trasformazione rigida $3\text{D}$ (isometria) composta da:
\begin{itemize}
    \item Una \textbf{Rotazione} espressa da una matrice $R \in \mathbb{R}^{3 \times 3}$.
    \item Una \textbf{Traslazione} espressa da un vettore $T \in \mathbb{R}^{3 \times 1}$.
\end{itemize}

I parametri di un sistema di acquisizione si dividono in due categorie:
\begin{itemize}
    \item \textbf{Parametri Estrinseci}: Definiscono la posizione e l'orientamento della telecamera nello spazio rispetto al sistema di riferimento Mondo ($R$ e $T$). Non dipendono dalla struttura interna della fotocamera.
    \item \textbf{Parametri Intrinseci}: Parametri propri della struttura interna del dispositivo (distanza focale $f$, coordinate del punto principale, fattori di scala dei pixel, eventuale sbieco/skew).
\end{itemize}

Il processo di determinazione sperimentale di questi parametri è detto \textbf{Calibrazione della Telecamera}.

\section{Punti di Fuga, Non-Linearità e Coordinate Omogenee}

\subsection{Non-Linearità e Punti di Fuga}
Nelle trasformazioni euclidee e affini le rette parallele rimangono parallele. Al contrario, le equazioni della proiezione prospettica dividono le coordinate spaziali per la profondità $Z$:
\begin{equation}
    (X, Y, Z) \longmapsto \left(f\frac{X}{Z}, \, f\frac{Y}{Z}\right)
\end{equation}
Questa divisione rende la proiezione prospettica una **trasformazione non lineare** nello spazio euclideo $\mathbb{R}^3$.

Di conseguenza, insiemi di rette parallele nello spazio $3\text{D}$ (non parallele al piano immagine) proiettate sull'immagine perdono il loro parallelismo e convergono in un punto comune chiamato **Punto di Fuga** (\textit{Vanishing Point}). Ogni direzione nello spazio $3\text{D}$ individua un punto di fuga distinto sul piano immagine.

\subsection{Coordinate Omogenee e Spazio Proiettivo}
Per poter trattare la proiezione prospettica in forma algebrica lineare e sfruttare il calcolo matriciale, si introducono le **Coordinate Omogenee** estendendo lo spazio euclideo allo spazio proiettivo ($\mathbb{P}^3$ e $\mathbb{P}^2$).

Un punto $3\text{D}$ di coordinate euclidee $(X, Y, Z)^T$ viene rappresentato in coordinate omogenee aggiungendo una quarta coordinata strumentale:
\begin{equation}
    \mathbf{\tilde{M}} = \begin{pmatrix} X \\ Y \\ Z \\ 1 \end{pmatrix} \sim \begin{pmatrix} \alpha X \\ \alpha Y \\ \alpha Z \\ \alpha \end{pmatrix}, \quad \alpha \neq 0
\end{equation}
Le coordinate omogenee definiscono classi di equivalenza a meno di un fattore di scala non nullo $\alpha$.

Grazie alle coordinate omogenee, la proiezione prospettica completa (comprensiva di parametri intrinseci ed estrinseci) si esprime mediante una semplice moltiplicazione matriciale lineare:
\begin{equation}
    \mathbf{\tilde{m}} = P \mathbf{\tilde{M}} = K [R \mid T] \mathbf{\tilde{M}}
\end{equation}
dove $P \in \mathbb{R}^{3 \times 4}$ è la matrice di proiezione prospettica e $K$ è la matrice dei parametri intrinseci.

\section{Effetti Reali della Lente e Imperfezioni Ottiche}

\subsection{Distorsione Radiale}
Le lenti reali (specialmente quelle economiche o di bassa qualità) introducono deformazioni geometriche note come **distorsione radiale**. La distorsione cresce allontanandosi dal centro dell'ottica verso i bordi esterni dell'immagine. Le coordinate misurate sul piano immagine reale deviano da quelle teoriche del modello pinhole e devono essere corrette analiticamente in fase di calibrazione.

\subsection{Modello della Lente Sottile e Messa a Fuoco}
In una fotocamera reale, l'apertura non è un punto infinitesimo ma possiede un'area finita. Per raccogliere luce sufficiente si utilizza una lente.

Nel modello teorico della **Lente Sottile**:
\begin{itemize}
    \item Un raggio parallelo all'asse ottico viene deviato passando per il fuoco.
    \item Un raggio che passa per il centro della lente prosegue privo di deviazione.
    \item Un raggio che attraversa il fuoco esce dalla lente in direzione parallela all'asse ottico.
\end{itemize}

Una lente fa convergere i raggi riflessi da un punto $P$ della scena su un preciso punto del sensore solo se $P$ si trova ad una specifica distanza (piano di messa a fuoco). Se il punto si trova più avanti o più indietro, i suoi raggi intercettano il piano del sensore formando un'area circolare detta **cerchio di confusione** (\textit{circle of confusion}), generando un effetto di \textbf{sfocatura} (\textit{blur}).

L'intervallo di distanze nello spazio entro il quale gli oggetti mantengono una sfocatura impercettibile ed appaiono nitidi è definito **Profondità di Campo** (\textit{Depth of Field}).

\section{Compromesso Fotografico: Apertura, Esposizione e Rumore ISO}
Nell'uso pratico della telecamera si presenta un compromesso fondamentale tra geometria ed esposizione luminosa:
\begin{enumerate}
    \item Per aumentare la profondità di campo (ad esempio per avere un panorama tutto perfettamente a fuoco), occorre **chiudere il diaframma** (ridurre il diametro dell'apertura).
    \item Riducendo l'apertura, la quantità di luce che attraversa l'ottica e raggiunge il sensore diminuisce drasticamente.
    \item Per compensare la scarsità di luce e ottenere un'immagine sufficientemente luminosa vi sono due soluzioni:
    \begin{itemize}
        \item \textbf{Aumentare la sensibilità ISO}: Si amplifica il segnale elettrico generato dal sensore. Tuttavia, poiché i sensori elettronici reali presentano imperfezioni e rumore di fondo, amplificare il segnale significa **amplificare anche il rumore**, producendo immagini sgranate, con colori alterati e degradati.
        \item \textbf{Aumentare il Tempo di Esposizione / Scatto}: Se la scena è statica, è preferibile prolungare la durata dell'acquisizione mantenendo i valori ISO al minimo. In questo modo il sensore accumula sufficiente energia luminosa garantendo un'immagine nitida, ben esposta e priva di rumore.
    \end{itemize}
\end{enumerate}
